\documentclass[conference,a4paper]{IEEEtran}
\usepackage[utf8]{inputenc}
\usepackage{amsmath,amssymb,amsfonts}
\usepackage{graphicx}
\usepackage{booktabs}
\usepackage{microtype}
\usepackage{hyperref}
\usepackage{cite}

\hypersetup{
    colorlinks=true,
    linkcolor=blue,
    citecolor=blue,
    urlcolor=blue
}

\begin{document}

\title{Dissecting and Adapting DINOv2 for Land Cover Classification: An Architectural Teardown and Parameter-Efficient Study on Sentinel-2 Imagery}

\author{
    \IEEEauthorblockN{Alfredo Jimenez, Jhonatan Davila, Italo Perez}
    \IEEEauthorblockA{
        \textit{Departamento de Ingenier\'{i}a, Pontificia Universidad Cat\'{o}lica del Per\'{u} (PUCP), Lima, Peru} \\
        \textit{Grupo de Investigaci\'{o}n en Inteligencia Artificial (IA-PUCP)} \\
        Corresponding email: italo.perez@pucp.edu.pe
    }
}

\maketitle

\begin{abstract}
Self-supervised Vision Transformers such as Meta AI's DINOv2 acquire robust visual representations on natural photography. However, deploying these representations to Earth observation imagery introduces a substantial domain shift due to nadir perspective, multi-scale agricultural textures, and subtle spectral class boundaries. In this work, we conduct an architectural teardown of DINOv2 (ViT-S/14) and examine a controlled nucleus intervention using Low-Rank Adaptation (LoRA) restricted to the deepest self-attention layers. Evaluated on the complete EuroSAT dataset (27,000 Sentinel-2 images) using an NVIDIA Tesla T4 GPU, our LoRA configuration adapts only 32,266 parameters (0.146\% of the backbone) while improving test Macro-F1 from 95.80\% to 97.14\% ($p = 0.002$) over a frozen linear probe, resolving fine-grained infrastructure confusion without inference overhead.
\end{abstract}

\section{Introduction}
Foundation vision models pretrained via self-supervised learning, particularly DINOv2~\cite{oquab2023dinov2}, yield versatile visual representations across natural image benchmarks. When transferred to Earth observation tasks, however, they encounter a significant domain shift: satellite imagery (e.g., Sentinel-2 EuroSAT~\cite{helber2019eurosat}) exhibits nadir viewing geometry, multi-scale spatial textures, the absence of focal foreground objects, and high intra-class spectral variance across land-cover classes. 

Standard linear probing (training an affine classifier over frozen features) prevents catastrophic forgetting but cannot adapt features to high-frequency remote sensing spatial distributions. Conversely, full fine-tuning modifies over 22 million parameters, risking overfitting on specialized target distributions and incurring heavy computational demands. In this work, we present an architectural teardown of DINOv2 (ViT-S/14) and evaluate a controlled nucleus intervention using Low-Rank Adaptation (LoRA)~\cite{hu2021lora} restricted to the deepest self-attention layers, demonstrating statistically significant performance gains with minimal parameter footprint and zero inference overhead.

\section{Methodology}
\subsection{Dataset and Protocol}
We benchmark on the complete EuroSAT RGB dataset~\cite{helber2019eurosat} (Sentinel-2, 10 Land Use and Land Cover classes, 27,000 images). Images are resized from $64 \times 64$ to $224 \times 224$ pixels via bicubic interpolation without cropping to preserve peripheral context, followed by ImageNet normalization. We adopt a deterministic stratified 70/15/15 split (seed 42): 18,900 images for training, 4,050 for validation, and 4,050 for testing. 

Following the canonical DINOv2 protocol~\cite{oquab2023dinov2}, the classifier input is formed by concatenating the class token with the spatial mean of patch tokens ($[\text{CLS} \,;\, \text{Mean}(\text{Tokens})]$, 768 dimensions), yielding an affine head of 7,690 parameters.

\subsection{Phase A: Architectural Teardown}
DINOv2-Small (ViT-S/14 without registers) comprises 22,056,576 parameters (88.2\,MB in fp32) organized across 12 transformer encoder blocks ($d=384$, 6 heads, MLP expansion $4\times$ to 1536). Analytical forward computation totals 12.25\,GFLOPs per 224\,px image (1.011\,GFLOPs per block, where quadratic attention represents 9.9\%). 

Empirical profiling on an NVIDIA Tesla T4 GPU yields a latency of 8.89\,ms (112.5\,img/s, batch 1, fp32) and 1.17\,ms per image (853.7\,img/s, batch 128, fp16). Effective Receptive Field (ERF) analysis~\cite{luo2016understanding} shows global coverage starting from block 1 ($r_{50}$ expands from 69.4\,px to 85.0\,px, and $r_{90}$ from 117.2\,px to 127.1\,px at block 12). Principal Component Analysis (PCA) reveals that while block 4 tokens encode spatial coordinates, block 12 tokens cluster into clear semantic class representations.

\subsection{Phase B: Nucleus Intervention via LoRA}
Hypothesizing that remote sensing domain adaptation requires realigning high-level semantic tokens rather than early geometric primitives, we freeze blocks 0--10 and inject low-rank decomposition matrices exclusively into the Query ($Q$) and Value ($V$) projections of blocks 11 and 12. Each adapted weight is parameterized as:
\begin{equation}
W = W_0 + \frac{\alpha}{r}(B \cdot A)
\end{equation}
where $W_0 \in \mathbb{R}^{384 \times 384}$ is frozen, $A \in \mathbb{R}^{8 \times 384}$ is Kaiming-uniform initialized, and $B \in \mathbb{R}^{384 \times 8}$ is zero-initialized ($r=8, \alpha=8$). This restricts trainable parameters to 32,266 (24,576 LoRA + 7,690 head, representing 0.146\% of the backbone). 

Both baseline and intervention are trained under identical conditions (AdamW, $\text{lr}=10^{-3}$, weight decay $10^{-4}$, batch size 64, 8 epochs) using precomputed frozen embeddings, accelerating extraction to 61 seconds.

\section{Results and Discussion}
\subsection{Quantitative Benchmarks}
As shown in Table~\ref{tab:results}, LoRA adaptation outperforms the frozen linear probe across all evaluation metrics on the 4,050 held-out test images. Test Macro-F1 increases from 95.80\% to 97.14\% ($+1.34\%$), and Test Accuracy improves from 96.00\% to 97.28\% ($+1.28\%$), reducing total test misclassifications from 162 to 110. A two-tailed McNemar test on paired test predictions confirms statistical significance ($\chi^2 = 9.56, p = 0.002$).

\begin{table}[t]
\centering
\caption{Performance on full EuroSAT test set (4,050 images, seed 42).}
\label{tab:results}
\begin{tabular}{lcccc}
\toprule
\textbf{Configuration} & \textbf{Params} & \textbf{\% Back.} & \textbf{Accuracy} & \textbf{Macro-F1} \\
\midrule
Linear Probe (Frozen) & 7,690 & 0.035\% & 96.00\% & 95.80\% \\
\textbf{LoRA (Bl. 11--12, $r=8$)} & \textbf{32,266} & \textbf{0.146\%} & \textbf{97.28\%} & \textbf{97.14\%} \\
\bottomrule
\end{tabular}
\end{table}

\subsection{Class-Level Dynamics and Error Analysis}
Disaggregating F1 performance by class indicates that the largest gains occur in complex linear infrastructure and hydrology: \textit{Highway} F1 increases by $+5.85\%$ (91.87\% to 97.72\%) and \textit{River} improves by $+3.71\%$ (93.80\% to 97.51\%). High-baseline natural categories (\textit{Forest} at 98.89\%, \textit{SeaLake} at 99.44\%) maintain peak fidelity without catastrophic forgetting. Residual confusion is predominantly concentrated among subtle spectral boundaries between \textit{Pasture} and \textit{HerbaceousVegetation}, where absence of short-wave infrared bands in RGB creates physical ambiguity.

\section{Conclusions}
Targeted Low-Rank Adaptation on the Query and Value projections of the final two transformer blocks bridges the domain gap between natural and remote sensing imagery, delivering a statistically significant gain ($+1.34\%$ Macro-F1, $p=0.002$) while tuning only 0.146\% of parameters. Embedding caching accelerates training by over $8\times$ without numerical deviation, and weight merging eliminates runtime latency penalties entirely.

\noindent \textbf{Limitations \& Phase C Roadmap:} We maintain strict epistemological honesty: the random stratified partition evaluates intra-distribution generalization. In geographic remote sensing, spatial autocorrelation between adjacent tiles can inflate empirical metrics. Ongoing work for Phase C investigates out-of-distribution robustness via geographically disjoint regional splits using EuroSAT 13-band Sentinel-2 L2A GeoTIFFs and cross-sensor evaluation on Sentinel-1 SAR data.

\bibliographystyle{IEEEtran}
\begin{thebibliography}{1}
\bibitem{oquab2023dinov2}
M.~Oquab et~al., ``DINOv2: Learning robust visual features without supervision,'' \emph{arXiv:2304.07193}, 2023.
\bibitem{hu2021lora}
E.~J. Hu et~al., ``LoRA: Low-rank adaptation of large language models,'' in \emph{Proc. ICLR}, 2022.
\bibitem{helber2019eurosat}
P.~Helber et~al., ``EuroSAT: A novel dataset and deep learning benchmark for land use and land cover classification,'' \emph{IEEE JSTARS}, vol.~12, no.~7, pp.~2217--2226, 2019.
\bibitem{luo2016understanding}
W.~Luo et~al., ``Understanding the effective receptive field in deep convolutional neural networks,'' in \emph{Proc. NeurIPS}, 2016.
\end{thebibliography}

\end{document}
